\chapter{Conclusion and Future Work}
\label{ch:conclusion}

\section{Conclusion}
\label{sec:conclusion_summary}

A real-time Cartesian impedance controller was implemented on a
seven-degree-of-freedom Franka Emika Panda to generate compliant
contact-induced rotation towards the configured surface. The controller assigns
directional translational and rotational stiffness and damping in the surface
frame and allows the virtual centre of compliance to be displaced from the TCP.

Shifting the \abbr{CoC} transforms the Cartesian stiffness and damping matrices
and introduces off-diagonal coupling terms between translation and rotation. A
translational position error can therefore contribute to the commanded moment,
while a rotational error can contribute to the commanded force. The shifted
reference also gives the commanded press the virtual lever arm
\mbox{\(r_c\)}, which generates the additional moment
\mbox{\(r_c\times f\)} at the TCP.

In the contact measurements, changing the \abbr{CoC} position had little
effect on the model-estimated normal force but changed the estimated TCP
moment substantially. The angular error
\mbox{\(\theta_{\mathrm{err},t_1}(t_{\mathrm{end}})\)} changed with this moment.
The comparison indicates that the \abbr{CoC} mainly influenced the angular
behaviour through its additional commanded moment at the TCP.

With the \abbr{CoC} at the TCP, the angular-error magnitude was smaller
than the entry magnitude by \mbox{\(81.2\,\%\)} and \mbox{\(85.0\,\%\)}
for the positive- and negative-offset conditions.
The negative-offset component crossed zero.
The nominal-zero condition instead showed an increase in angular error.
At \mbox{\(r_c=0\)}, the additional CoC-induced moment vanishes while the
rotational spring and damper remain active. The TCP provides a neutral
impedance reference with respect to this added moment.

Higher rotational stiffness left a larger angular error.
Increasing \mbox{\(K_{R,t_1}\)} from \mbox{\(5\)} to
\mbox{\(50\,\mathrm{N\,m/rad}\)} increased the mean angular error by
\mbox{\(276.2\,\%\)} relative to the lowest stiffness.
The tested variation of translational stiffness along \mbox{\(t_2\)}
produced a span of \mbox{\(3.8\,\%\)} of its TCP-centred reference error.
The translational stiffness along \mbox{\(t_1\)} remained constant.
Within these tested ranges, the angular error was more sensitive to
rotational stiffness.

The effect of \abbr{CoC} displacement depended on the entry direction.
The smallest measured mean error magnitudes were approximately
\mbox{\(15\,\%\)} and \mbox{\(34\,\%\)} below their respective TCP-centred
references, at different tested positions.
The first-tangent error component could cross zero.
Further contact-induced rotation could increase its final magnitude.
The illustrated positive-offset trace reached a band around its angular error
earlier with a selected displacement than at the TCP.
These results motivate selection based on the measured angular error and its
evolution during contact. Returning the \abbr{CoC} to the TCP after alignment
would remove the directional preference of the added moment.

The secondary null-space functions were evaluated separately with a fixed
Cartesian pose reference. Projected damping reduced the cumulative projected
null-space motion \mbox{\(E_N\)} by \mbox{\(25.1\,\%\)}, while the net
displacement decreased by \mbox{\(25.4\,\%\)}. The cumulative motion
and net displacement were nearly equal, meaning that only a small part of the
measured motion was cancelled by motion in the opposite direction. Projected
damping therefore reduced the null-space motion generated by the disturbance.

Singular-value conditioning was active before and during the virtual
disturbance. Both settings left mean net displacements \mbox{\(\Delta\eta\)}
close to zero under the tested direction. At
\mbox{\(k_\sigma=1.5\,\mathrm{N\,m}\)}, the value was
\mbox{\(0.015^\circ\)}. At \mbox{\(k_\sigma=2.0\,\mathrm{N\,m}\)}, it was
\mbox{\(-0.011^\circ\)}. The higher conditioning torque gave the smaller
magnitude of the mean net displacement. However, the cumulative projected
null-space motion \mbox{\(E_N\)} increased by \mbox{\(486.6\,\%\)}.
This is consistent with greater back-and-forth motion while the mean net
displacement remained close to zero. The minimum singular value
\mbox{\(\sigma_{\min}\)} remained close to its disturbance-entry value
throughout the disturbance in both settings.

\section{Limitations}
\label{sec:limitations}

\subsection{Scope of the Surface-Contact Evaluation}
\label{subsec:conclusion_experimental_range}

The surface-contact study used one Panda configuration, one rectangular tool
and mounting arrangement, one configured surface reference, and one Contact
Establishment trajectory. The measured effects therefore apply to this
experimental geometry and to the investigated ranges of rotational stiffness,
translational stiffness, angular offset, and \abbr{CoC} displacement.

The main repeatable comparison concerned rotation about \mbox{\(t_1\)}.
Measurements about \mbox{\(t_2\)} showed greater variability and did not
provide the same repeatable response. Each contact trial used one selected
\abbr{CoC} displacement throughout Contact Establishment; an adaptive change of
the \abbr{CoC} during alignment was not evaluated. The quantitative evaluation
focused on Contact Establishment and did not include the subsequent grinding
motion.

\subsection{Calibration and Measurement Scope}
\label{subsec:conclusion_measurement_uncertainty}

The measured angular offset \mbox{\(\theta_{\mathrm{meas},t_1}\)} and angular
error \mbox{\(\theta_{\mathrm{err},t_1}(t)\)} use the calibrated surface
reference and tool normal with the measured end-effector orientation.
The surface reference was established with the tool face seated on the
physical plate and provides an estimate of its normal.
Its construction used the nominal end-effector axis, while the separate
tool-normal calibration found a direction differing by approximately
\mbox{\(1.56^\circ\)}.
The configured reference retained its original construction.

The instantaneous physical tool--surface angle was not measured independently
during contact. The calculated angular error therefore estimates physical
misalignment subject to the surface-reference and tool-normal calibration.
No numerical uncertainty bound for that estimate was established.
The reported component about \mbox{\(t_1\)} also leaves the error about the
other tangent outside the principal comparison.
A zero first-tangent component alone does not establish complete normal
alignment.

The comparison endpoints were stored to \mbox{\(0.01^\circ\)}.
A sample standard deviation of zero at this precision indicates unresolved
variation in the stored endpoint values.
The physical pressure centre and instantaneous contact location were not
measured, so their individual contributions to the interaction remain
unresolved.

\subsection{Tool-Mount Compliance}
\label{subsec:conclusion_tool_mount}

The tool mount exhibited approximately \mbox{\(\pm2^\circ\)} of unloaded
angular play about \mbox{\(y_{\mathrm{EE}}\)}, which is approximately aligned
with \mbox{\(t_2\)} in the flat configuration. This mechanical play may have
contributed to the greater variability of the \mbox{\(t_2\)} measurements.

The angular error \mbox{\(\theta_{\mathrm{err},t_1}(t)\)} assumes the
calibrated tool normal remains fixed relative to the end effector.
The instantaneous relative tool--gripper rotation was not tracked separately
during contact. Motion within the mount can therefore make the physical
tool-face direction differ from the direction calculated using that calibration.

\subsection{Scope of the Null-Space Evaluation}
\label{subsec:conclusion_nullspace_scope}

Projected damping and singular-value conditioning were evaluated separately in
Cartesian pose hold under one virtual disturbance direction and one robot
configuration. The measured results therefore describe the controller behaviour
for the investigated condition rather than under arbitrary disturbances.

\section{Future Work}
\label{sec:future}

For future work, a more rigid tool mounting arrangement would provide
a more stable basis for the evaluation about \mbox{\(t_2\)}, where greater
variability was observed. The \abbr{CoC} results further motivate an adaptive
strategy in which the controller begins with the \abbr{CoC} at the TCP,
determines the angular-offset direction, and temporarily applies a tangential
\abbr{CoC} displacement that supports the required rotation. After alignment,
the \abbr{CoC} can return to the TCP to remove the directional preference and
retain compliant behaviour for both rotational directions.
